\section{Willingness to contribute: majority willing to contribute a few percent}

The preceding section showed that support for international climate policy can persist when its costs are made explicit. This section moves from acceptance of a collective policy to stated willingness to bear a personal contribution. The two concepts should not be conflated: policy acceptance measures whether respondents endorse a collective rule, whereas willingness to contribute captures the personal sacrifice they report being prepared to make. 
% The incidence of that sacrifice is equally important. % AF: deletion: unclear
% A distributionally fair scheme should protect low-income households and concentrate its net burden on those with greater ability to pay rather than treating a uniform contribution as distributionally neutral.

Using a globally representative survey of nearly 130,000 respondents across 125 countries, \citet{andre_globally_2024} ask whether individuals would contribute 1\% of their personal income to climate action. Globally, 69\% report that they would contribute, although country-level willingness ranges from 30\% to 93\%, the result is therefore majoritarian worldwide but not in every country. % AF: dire explicitement quels pays sont sous 50% et à combien ils sont (et dire qu'ils passent > 50% lorsque le coût est <1%)
The most consequential finding concerns second-order beliefs: respondents estimate that only 43\% of others would contribute, creating a 26-percentage-point gap between reported and perceived willingness. Barriers to collective climate action may therefore arise not only from private costs but also from systematic underestimation of the social support on which cooperation depends. Consistent with this interpretation, \citet{mildenberger_beliefs_2017} document widespread underestimation of pro-climate beliefs among mass publics and elites in the United States and China and show experimentally that correcting second-order beliefs increases support for climate-policy action.
% TODO: On WTP, World Bank (09) has similar results (p. 16). https://documents1.worldbank.org/curated/en/492141468160181277/pdf/520660WP0Publi1und0report101PUBLIC1.pdf (also interesting things p. 5, 10, 11, 13, 16)
% Results in Figures 1, 2 and 3 and Table S3. 11 countries don’t have majority WTC, incl. the UK (48%), the U.S. (48%), Canada (49%), Japan (49%), Israel (37%), Pakistan, New Zealand, Kazakhstan, Russia. Pluralistic ignorance is ~19 pp in U.S. and Eu4, i.e. similar to what we find for the Eu4 in Fabre, Douenne, Mattauch (but larger for the U.S.).
% “Would you be willing to contribute 1% of your household income every month to fight global warming? This would mean that you would contribute [$1] for every [$100] of this income.” Yes/No/PNR

Information does not, however, affect every form of international redistribution uniformly. In the United States, \citet{nair_misperceptions_2018} finds that information about the global income distribution increases support for foreign aid and raises real-stakes donations to foreign charities by 55\% relative to the control group. By contrast, \citet{fehr_your_2019} show that Germans systematically underestimate their position in the global income distribution and express a positive willingness to pay for information about their national and global income ranks. Experimentally correcting global-rank misperceptions nevertheless affects neither support for policies addressing global inequality nor giving to the global poor. The demand for information documented by \citet{fehr_your_2019} should therefore not be interpreted as evidence of a greater willingness to contribute. Taken together, these studies indicate that information can activate international distributive concern in some settings, but that correcting perceived global rank is not a general mechanism for increasing support for global redistribution.

Willingness to contribute remains cost-sensitive, but cost is only one determinant of policy acceptance. \citet{bechtel_mass_2013} show that higher household costs reduce support for global climate agreements, whereas equitable burden-sharing, broad participation, and credible enforcement increase it. % AF: give numbers
At the same time, support for unilateral climate action remains substantial when costs and free-riding are made explicit and appears to depend more strongly on expected effectiveness and domestic co-benefits than on strict reciprocity \citep{bernauer_how_2015,mcgrath_how_2017}. % AF: detail how they arrive at this conclusion
Support for North--South climate finance is similarly higher when recipient governments are perceived as effective, allocation decisions are jointly governed, and other donor countries contribute substantial resources \citep{gampfer_obtaining_2014}. Willingness to contribute is therefore neither unconditional nor reducible to price; it depends on policy incidence, expected effectiveness, institutional credibility, and beliefs about collective participation.

Perceptions further condition willingness to contribute internationally. Citizens tend to overestimate the foreign aid already provided by their country \citep{kaufmann_foreign_2019}, while individuals in high-income countries frequently underestimate their position in the global income distribution \citep{nair_misperceptions_2018,fehr_your_2019}. The null treatment effect documented by \citet{fehr_your_2019} nevertheless shows that correcting the latter misperception is insufficient, by itself, to increase support for global redistribution. More broadly, \citet{fabre_public_2026} finds that global climate and redistributive policies can attract majority acceptance even in countries whose citizens are expected to be net contributors. Across this literature, acceptance is stronger when burdens are distributed fairly, policies are expected to be effective, and implementation arrangements are institutionally credible \citep{fabre_public_2026,bechtel_mass_2013,gampfer_obtaining_2014}. The evidence therefore points to a conditional willingness to contribute rather than to unconditional or politically intense demand. Moreover, this willingness may remain politically latent when citizens underestimate its prevalence or attach insufficient electoral salience to international redistribution.

